\section{Evaluation}
\label{sec:evaluation}

\needsdata{this entire section. The structure below is the plan; the tables
are wired to the artefact generator and currently contain a smoke-test run
against the local reference adapter, which is not a provider comparison.}

We evaluate three hosted providers --- E2B, Daytona and Runta --- plus the
local reference adapter as a latency floor.

\subsection*{RQ1: Do serverless metrics predict sandbox performance?}

This is the question a reviewer asks first: why another suite? We test it
directly. For each provider we measure a FaaS-style baseline (cold start plus
one single-shot execution) and three sandbox-shaped workloads (100 sequential
short commands; install, edit, resume; reset followed by 100 steps), and ask
whether the provider ranking is stable across them and whether cold-start
latency correlates with session completion time.

If the rankings agree and the correlation is strong, the honest conclusion is
that existing metrics suffice and this suite is not needed. If they diverge,
SandboxBench has an experimental justification rather than an argued one.
\needsdata{which of the two it is}

\subsection*{RQ2: How do bottlenecks move with session shape?}

We sweep command count, interaction frequency and state size, and locate the
crossover where per-command overhead overtakes create latency in the
user-visible total.

\subsection*{RQ3: How does single-sandbox performance relate to cluster throughput?}

We compare each provider's single-sandbox dimension results against its burst
and sustained provisioning results, and test whether the fastest single
sandbox belongs to the provider with the highest sustainable throughput.

\subsection*{RQ4: How do lifecycle strategies affect rollout throughput and cost?}

We run the reinforcement-learning profile under all four lifecycle profiles
and compare transitions per second and cost per valid rollout across the three
reset mechanisms.

\subsection{Results}

\begin{table}[t]
\caption{Headline results. Generated by \texttt{sandboxbench report
--format=latex}; do not edit by hand. \needsdata{a real cross-provider run ---
the current contents are a local smoke test}}
\label{tab:results}
\small
% Placeholder. No benchmark run has produced this table yet.
%
% Regenerate it from a real run:
%   make -C paper table RESULT=results/<run>.json
%
% Do not hand-write numbers here. The point of generating it is that no value
% is retyped between the artefact and the paper.
\begin{tabular}{lrrr}
\toprule
Benchmark & E2B & Daytona & Runta \\
\midrule
\textit{no run yet} & --- & --- & --- \\
\bottomrule
\end{tabular}

\end{table}

\subsection{Threats to validity}

\begin{description}
  \item[Provider configuration is not controlled.] Hosted providers do not
  expose CPU generation, storage class or prewarming policy. A difference
  between two providers may be a hardware difference.
  \item[The client is a confounder.] Every number includes a public network
  path from one client location. We report the location and repeat from more
  than one. \needsdata{how many}
  \item[Synthetic workloads may not be representative.] Where
  \S\ref{sec:background:characterisation} could not be answered from real
  traces, the corresponding workload is synthetic and labelled as such.
  \item[Measurements age.] These are commercial services under active
  development. A result is a measurement of a date, and the date is in every
  result file.
\end{description}
